\customchapter{RESUMEN}

La documentación digital del patrimonio construido depende de levantamientos
láser cuya interpretación sigue siendo manual: el escáner registra dónde hay
materia, pero no qué elemento representa cada punto. Los métodos de
segmentación semántica tridimensional usan una lista cerrada de categorías y
asignan todo elemento imprevisto a la clase más parecida. El objetivo es
mejorar de forma medible el reconocimiento de los elementos sin etiqueta.

La investigación es aplicada, cuantitativa y de diseño experimental
comparativo. HyperNCD se entrenó con SemanticPOSS y, al comprobar su
insuficiente transferencia desde escenas urbanas, con ArCH, un conjunto de
patrimonio arquitectónico~\cite{matrone2020benchmark}. El modelo semántico y el
codificador auto-supervisado Sonata se ejecutaron congelados y sus descriptores
se fusionaron después, sin entrenamiento conjunto.

En una sección de la Catedral de Lima, el checkpoint de SemanticPOSS asignó
98,4\,\% de los puntos a una sola cabeza; el de ArCH (época 30) la redujo a
44,6\,\%, elevó el NMI sin clase residual de 0,027 a 0,439 y alcanzó 35,5\,\% de
mIoU uno a uno. En diez semillas de agrupamiento, M2 obtuvo $48,4\pm1,1$\,\%, M3
$45,4\pm8,5$\,\%, M4 base $51,4\pm1,7$\,\% y la fusión M4 $52,5\pm2,5$\,\%. La
diferencia de 1,1 puntos entre M4 y M4 base cae dentro de la dispersión y no es
una mejora concluyente.

La adaptación de dominio con ArCH evita el colapso observado con SemanticPOSS y
hace viable evaluar el sistema híbrido, con alcance limitado a una escena, una
corrida de entrenamiento y repeticiones solo del agrupamiento.

\vspace{0.5em}

\noindent\begin{minipage}{\linewidth}
\noindent\textbf{Palabras clave:}\par
\noindent Nubes de puntos; Segmentación semántica tridimensional; Descubrimiento
de clases novedosas; Aprendizaje auto-supervisado; Hipergrafos; Patrimonio
arquitectónico
\end{minipage}
